\section{Future Work}
\label{sec:future}

We order this by what would most change the paper's claims rather than by
implementation convenience.

\para{Execute the evaluation.} The protocol in \appref{app:protocol} is the
immediate obligation. E3(i) --- whether authored judges generalize off their
calibration corpus --- is the measurement we most want, because a strong result
there would undercut this paper's own argument for keeping a human at \mode{Apply}
(D4). Designing an experiment whose outcome could weaken your position is the only
way the position ends up meaning anything.

\para{Multi-environment promotion.} The highest-value architectural work, and the
most involved. A \code{dev}\arrowto\code{staging}\arrowto\code{prod} ladder needs
per-stage configuration and secrets, approval routing with a requester-and-approver
distinction, promotion state that survives stage transitions, and cross-stage
compatibility tests. Because the current prompt record is a born-approved provenance
anchor, this is a rebuild of the approval path rather than the enabling of a flag
(D9). The interesting design question is whether stages should be a first-class
dimension of the \livetarget{} or a separate axis --- the former is simpler and the
latter composes better with per-family release idioms.

\para{Capability interfaces over the asset families.} \secref{sec:arch:claim}
identifies the design \sys should have had: a thin base contract every \gasset{}
satisfies, plus orthogonal capability interfaces --- \code{Releasable},
\code{Rollbackable}, \code{EvidenceBearing} --- that a family declares and is then
\emph{required} to implement. The payoff is not elegance. It converts
\tabref{tab:families} from documentation a reader must consult into a constraint a
family cannot violate, which is the difference between a guarantee surface that is
described and one that is enforced. The work is a retrofit across
\statFamilies{} families and their routes, and it subsumes the registry projection
described next.

\para{Complete the capability registry.} Converging the \statAriaTools{} hand-written
tools onto the registry projection would make the seam real and, more importantly,
would let risk-tier classification be checked mechanically. That directly addresses
the classification gap in \secref{sec:limitations}: a projection with mandatory tier
declaration makes an unclassified mutation impossible to introduce, which is a
stronger property than any amount of review.

\para{Isolation as a first-class backend.} The sandbox protocol already exists; what
is missing is a shipped OS-enforced implementation. A microVM or seccomp-based
backend~\cite{agache2020firecracker,bubblewrap} would let \sys offer isolation rather
than containment for untrusted tool authorship, which is currently the platform's
sharpest security limitation.

\para{Attestable provenance.} Provenance anchors are internal records. Signing them
and emitting in-toto-style attestations would make an artifact's history verifiable
by a party that does not trust the \sys instance that produced
it~\cite{torresarias2019intoto,slsa}. This matters most for organizations that need
to demonstrate governance to an external auditor, which is the population most likely
to want a system like this.

\para{Cross-family impact analysis.} Today a candidate is evaluated against its own
evidence. A prompt change, a knowledge-base rebuild, and a tool revision can
interact --- and no current mechanism scores the \emph{bundle}. This is the
architectural precondition for multi-agent bundle optimization, and it is also the
work most likely to require revisiting the per-family factoring, since a bundle gate
would be the first genuinely cross-family enforcement point.

\para{Global rate limiting.} Moving the in-memory limiters to a shared store would
remove the per-process ceiling described in \secref{sec:exec}. This is
straightforward, and we list it because leaving a known sharp edge documented rather
than fixed is a legitimate criticism of the current implementation.

\para{Broader adversarial evaluation.} Injection through retrieved and tool-returned
content, malicious tool authorship, and compromised MCP servers all warrant
measurement rather than the architectural argument \secref{sec:components} currently
offers~\cite{greshake2023injection}. The interesting question is not whether \sys
prevents injection --- it does not --- but how much the capability model actually
bounds the blast radius, which is a quantity nobody has measured for any comparable
system.

\para{Automating \mode{Apply}, conditionally.} We would revisit D4 only under two
conditions, and we state them so the position is falsifiable rather than
ideological: calibrated per-family confidence that accounts for corpus coverage, and
a demonstrated off-corpus generalization result from E3(i). Absent both, an
auto-apply threshold trades a reviewable change for an invisible one, which is the
opposite of what this system is for.
